\cotefile{136}
\scanpage{235}

On considère l'anneau ½simplicial

\[
 \widehat A_*^*(S,\gamma^*,t)\deleted{\unsure{=|\mathrm{dfn}}}=\widehat A_*^*
\]

par

\[
\begin{aligned}
\widehat A_n^* &\simeq S\{\!\{X_0,\ldots,X_n\}\!\}\big/\bigl(\text{idéal \inserted{à p.d.,} \deleted{\unsure{fermé engendré}} engendré par }\textstyle\sum X_i-T\bigr)
 \otimes_S\frac{(\bigwedge^*[dX_0,\ldots,dX_n)}{\text{idéal }(\sum dX_i)}\\
 &(\simeq S\{\!\{Y_0,\ldots,Y_{n\unsure{-1}}\}\!\}[dY_0,\ldots,dY_{n-1}]
\end{aligned}
\]

On a donc \deleted{$\simeq$}

\[
 \widehat A_*^*\simeq\widehat A_*^0\otimes_S\bigwedge^*\Psi_*^S
\]

d'où

\[
 \pi_\ell(\widehat A_*^i)=\pi_\ell(\widehat A_*^0\otimes\bigwedge^i\Psi_*^S)\simeq\pi_{\ell-i}(\widehat A_*^0)
\]

et il faut donc étudier $\pi_*(\widehat A_*^0)$.

On considère sur $\widehat A_*^0$ \deleted{$\unsure{\Gamma_S^*(\Phi_*)}$} la filtration décroissante par idéaux fermés stables par p.d. ``engendrée'' par la condition que \deleted{\unsure{pour}} pour \unsure{tt|tout} $n$, les $X_i$ et $t$ soient en filtration $\geq 1$. On

\[
\begin{aligned}
\operatorname{Fil}^0(\widehat A_n^0)&\simeq \deleted{\unsure{S}}\widehat A_n^0\\
\operatorname{Fil}^p(\widehat A_n^0)&\simeq\left\{\sum\deleted{\unsure{a}}a_{i_0\ldots i_n}\deleted{\unsure{\{}}X_0^{(i_0)}\ldots X_n^{(i_n)}\ \middle|\ a_{i_0\ldots i_{\unsure{p|n}}}\in S^{[p-\sum_0^n i_\alpha]}\right\}\\
 &=\left\{\sum a_{i_0,\ldots i_{n-1}}Y_0^{(i_0)}\ldots Y_{n-1}^{(i_{n-1})}\ \middle|\ a_{i_0\ldots i_{\unsure{p|n|n-1}}}\in S^{p-\sum_0^n i_\alpha}\right\}
\end{aligned}
\]

Bien sûr, la topologie \unsure{produit|qui provient de} sur $\widehat A_n^{\unsure{*,0|*}}$ est aussi celle définie par cette filtration, et par suite $\widehat A_n^0$ est séparé et complet pour cette filtration. Un calcul immédiat donne, posant

\[
 k=S/S^{[1]}
\]

\[
 \operatorname{Gr}^{\unsure{p|1}}(\widehat A_*^0)
 \inserted{\simeq\widehat A_*^{p0}(\underbrace{\operatorname{Gr}^*(S),\widetilde\gamma^*,\widetilde t}_{\substack{k\text{-PD-algèbre}\\\text{\uline{graduée}}\\\text{quotient de }k\{T\}}})}
 \simeq\Gamma_k^p(\underline\Phi_*^k)
\]

\inserted{si l'épi $k\{T\}\longrightarrow\operatorname{Gr}^*(S)$ est un \uuline{iso}} \deleted{\unsure{d'où}}

d'où, dans cette hypothèse \inserted{d'isomorphisme}

\[
 \pi_i(\operatorname{Gr}^p(\widehat A_*^0))=
 \begin{cases}0&\text{si }p\ne0,\ i\ne0\ \deleted{\unsure{ou}}\\
 k&\text{si }(i,p)=0
 \end{cases}
 \quad(\text{donc }\operatorname{Gr}^p(\widehat A_*^0)\text{ flasque si }p\ne0)
\]

d'où on conclut que \deleted{l'inclusion}

\[
 \deleted{\operatorname{Fil}^1\widehat A_*^0\subset\widehat A_*^0}
\]

% archive sheet number: 234
